\section{Introduction}\label{sec:introduction}

A \emph{projection game} consists of a finite bipartite multigraph with disjoint
vertex sets $U,V$, a nonempty multiset $E$ of edge occurrences, finite answer
sets at its vertices, and a map $\pi_e$ from the left answer set to the right
answer set for each occurrence $e=(u,v)$.
A labeling chooses one answer at each vertex and satisfies $e$ when the right
answer equals the image of the left answer under $\pi_e$.
The \emph{value} of the game is the maximum fraction of satisfied occurrences.
Parallel occurrences may have different projection maps.

A \emph{2-to-1 game} has a common left alphabet $[2q]$ and right alphabet $[q]$,
where $[m]=\{1,\ldots,m\}$, and
\[
             |\pi_e^{-1}(b)|=2
        \qquad(e\in E,\ b\in[q]).
\]
Khot introduced the games conjectures in the study of hardness of
approximation~\cite{Khot2002}.  His original formulation allowed fibers
of size at most two; we use the exact-two convention.  The 2-to-1 Games
Conjecture asks for
NP-hardness of distinguishing satisfiable games from games of arbitrarily
small fixed value.  The satisfiable case is the \emph{perfect-completeness}
requirement: one labeling must satisfy every constraint.  Completeness
$1-\eps$ for each fixed $\eps>0$ does not itself give this endpoint, since
those reductions need not produce any finite satisfiable game.
There is also an obstruction to amplifying a constant gap by ordinary
repetition: repeating an exact-two projection $k$ times gives $2^k$
preimages of every right answer.  We prove the conjecture while retaining
two-element fibers, in the following explicit form.

\begin{theorem}\label{thm:main}
For every fixed rational $\delta\in(0,1)$ there are an integer $q=q(\delta)\ge2$
and a deterministic polynomial-time reduction from $3$-SAT to 2-to-1 games
with alphabets $[2q]$ and $[q]$ such that:
\begin{enumerate}[label=\textup{(\roman*)}]
\item a satisfiable formula produces a game of value $1$;
\item an unsatisfiable formula produces a game of value at most $\delta$.
\end{enumerate}
The output contains a nonempty explicit list of unweighted edge occurrences
and an explicit table for every projection map.
Every table has exactly two preimages for each right answer.
The polynomial running-time bound is in the ordinary binary input length;
its degree and constants may depend on the fixed $\delta$.
\end{theorem}

The theorem supplies the hypothesis of several perfect-completeness
reductions.  We derive consequences for maximum $k$-colorable subgraph
and satisfiable Not-Two constraints below, identifying those already
known unconditionally.

\subsection{Historical context and ingredients}

Austrin, O'Donnell, Tan and Wright established perfect-completeness 2-to-1
Label Cover hardness with alphabets of sizes $6$ and $3$ and soundness
$23/24+\eps$~\cite[Theorem~1.3, author manuscript]{AustrinODonnellTanWright2014}.
The remaining issue was to make soundness arbitrarily small while retaining
both perfect completeness and two-element fibers.

The near-perfect-completeness theorem arose from the Grassmann and shortcode
program.  Khot, Minzer and Safra introduced a Grassmann-based candidate
reduction~\cite{KhotMinzerSafra2017}; Dinur, Khot, Kindler, Minzer and Safra
developed the reduction to a Grassmann agreement hypothesis and the
associated expansion problem
\cite{DinurKhotKindlerMinzerSafra2018,DinurKhotKindlerMinzerSafra2021}.
Barak, Kothari and Steurer related the agreement problem to expansion through
the shortcode viewpoint~\cite{BarakKothariSteurer2018}.
The near-perfect expansion theorem of Khot, Minzer and Safra completed these
ingredients, giving completeness arbitrarily close to one and arbitrarily
small soundness~\cite{KhotMinzerSafra2018}.
More recently, Fei, Minzer and Wang proved perfect-completeness hardness for
4-to-1 games at arbitrarily small fixed soundness
\cite[Theorem~1.6]{FeiMinzerWang2026}.

Our conclusion concerns ordinary 2-to-1 games.  In the Rich 2-to-1
framework of Braverman, Khot and Minzer, a uniformly random edge incident
to each fixed left vertex must induce a partition that is uniform among
all partitions of the left alphabet into pairs
\cite[Definition~5]{BravermanKhotMinzer2021}.  No such condition is part
of Theorem~\ref{thm:main}.  For any positive real soundness target, one may apply
the theorem with a smaller fixed rational target.

Our external theorem inputs are a perfect-completeness PCP gap,
projection-game parallel repetition, and Grassmann expansion
\cite{Dinur2007,DinurSteurer2014,KhotMinzerSafra2018}.
From the last input we derive the inverse shortcode estimates using the
chart and averaging arguments of Barak, Kothari and Steurer
\cite{BarakKothariSteurer2018}; Section~\ref{sec:preliminaries}
gives the forms needed here.
The sampling comparisons have antecedents in the sub-code covering
framework of Khot and Safra~\cite{KhotSafra2013}.  The passage from
agreement of restricted bilinear forms to a fixed form of bounded rank
has a close counterpart in Fei, Minzer and Wang
\cite[Lemmas~5.15--5.16]{FeiMinzerWang2026}.
We prove the comparisons needed for the function spaces and sampling
laws below.

\subsection{Proof overview}

We start with a clause-versus-variable game: the left player answers a
clause with a satisfying assignment to its variables, the right player
answers a variable with a bit, and the constraint checks their agreement
on that variable.  A perfect-completeness PCP and parallel repetition
make the source value a fixed small constant on NO inputs.  We place
independent copies of its questions at the leaves of a finite rooted tree.
Each leaf carries all binary functions on its local answer domain.
At an internal node we take the sum of the spaces spanned by products of
pairs of functions from each child.  In particular, the product of two functions from one descendant space
belongs to every space above that descendant.

At every nonleaf node we display an array, meaning a list of functions
from its space.  A question is the unordered partition of the product of
all leaf-answer domains according to the joint evaluation of these arrays.
A legal answer is a nonempty part, so one assignment realizes all its
restored outputs simultaneously.  We discard input coordinates that these
evaluations cannot distinguish; a clause coordinate seen only through one
variable is recorded by that variable.  This identifies a question on
projected source domains with its pullback to the original domains.
Deleting one binary output direction merges at most two parts and gives
the preliminary at-most-two constraints.  A satisfying source assignment
answers every constraint.  The same joint partitions also remain unchanged
when an upper row is shifted by products of displayed lower rows, forcing
exact identities on all restored responses.  Section~\ref{sec:construction}
defines the questions, their sampling law, and these identities.

To analyze soundness, we sample a random root-to-leaf path.  Arrays off
that path are uniform; arrays on it use a recursive sampler with fresh
randomness for every call.  Stopping the calls at a node expresses every
ancestor array as a deterministic function of a bounded number of draws
there.  Increasing the number of children hides the special path child
and makes the observed draws close to uniform.  A separate comparison
replaces the source questions below a small random set of children by
their right questions.  After lifting the resulting functions to the
original domains, this change also has small total variation.
Section~\ref{sec:sampling} proves these comparisons and bounds the chance
that a child selected from the observed functions was projected.

Suppose a labeling satisfies a fixed positive fraction of the preliminary
constraints.  With enough levels, some pair of levels has positive
simultaneous success.  At the upper level an inverse shortcode theorem
finds one response occurring with positive density on an affine slice.
The table records responses at every lower-level descendant of the upper
node, so its range may be large;
Theorem~\ref{thm:inverse} has constants independent of that range and also
applies to small quotient spaces.  Joint partition identities convert the
slice into a coset of linear functionals $z$ on the upper function space
with $z(1)=1$.  Each functional
defines a product form $F(g,h)=z(gh)$ on the lower function space.
For the displayed lower rows $s_1,\ldots,s_\ell$ and their restored
response $y$, joint legality gives
\[
                   \bigl(F(s_t,s_u)\bigr)_{t,u}=yy^{\mathsf t},
\]
so this tested matrix has rank at most one.  A Fourier decoder on the
projected side selects from a bounded list computed under a small-bias
law determined by its own local input (Lemma~\ref{lem:heavy-list}).
The two local forms agree after projection with fixed positive probability.
Section~\ref{sec:upper} constructs these decoders, with an agreement bound
independent of the lower row count.

The remaining difficulty is that the left form depends on the very rows
on which its rank is tested.  Passing that test therefore does not
immediately bound its full rank.  Resampling within the lower omitted-direction
fiber gives equal restrictions with positive probability.  The sparse-projection
posterior bound promotes this to equality of full forms, apart from a small
error.  A second inverse application then finds one output form occurring densely
on an affine slice of row lists.  A fixed high-rank form almost never has a tested
matrix of rank at most one on that slice.  Thus a positive part of the
agreement must come from bounded-rank forms (Section~\ref{sec:lower}).

A low-rank normalized product form selects a short odd list of answers at
each source leaf.  On a projected leaf, agreement gives the parity
pushforward of that list: retain exactly those right answers having an
odd number of preimages in the left list.  Some projected children have
all their sampled contributions equal to zero.  After exposing the other
data, both decoders can reconstruct their original inputs from their own
questions at designated leaves in these children.  The designated source
pairs remain independent because the probability of the zero event is
exactly independent of them (Lemmas~\ref{lem:clean-law}
and~\ref{lem:reconstruction}).  Projection-game repetition then contradicts
positive agreement (Section~\ref{sec:clean}).

Finally, Section~\ref{sec:assembly} chooses the constants in their dependency
order, completes at-most-two maps to exact 2-to-1 maps, and replaces rational
edge weights by an explicit unweighted multiset.

\subsection{Applications}

\begin{corollary}[Maximum $k$-colorable subgraph]\label{cor:max-k-colorable}
There are absolute constants $C>0$ and an integer $k_0\ge3$ such that, for
every fixed integer $k\ge k_0$, it is NP-hard to distinguish the following
cases for a finite unweighted graph $G=(W,F)$ with $F\ne\varnothing$:
\begin{enumerate}[label=\textup{(\roman*)}]
\item $G$ is $k$-colorable;
\item every coloring $\chi:W\to[k]$ satisfies
\[
 \frac{|\{\{u,v\}\in F:\chi(u)\ne\chi(v)\}|}{|F|}
 \le 1-\frac1k+C\frac{\log k}{k^2}.
\]
\end{enumerate}
The fraction on the left is the proportion of edges properly colored by
$\chi$, or equivalently cut by its partition into $k$ color classes.
\end{corollary}
\begin{proof}
Theorem~\ref{thm:main} supplies the exact-two, perfect-completeness premise
of Guruswami and Sinop's reduction
\cite[Conjecture~3.4 and Theorem~3.26]{GuruswamiSinop2013}.
Their regularization and reduction preserve perfect completeness, giving
$k$-colorable YES graphs, and their Remark~2.6 transfers the hardness to
unweighted graphs. Their soundness bound
$1-1/k+O(\log k/k^2)$ is therefore unconditional. Choose an absolute
upper bound $C$ for the lower-order term and enlarge $k_0$ so the two
cases are disjoint. The reduction parameters, including the soundness
used in Theorem~\ref{thm:main}, are fixed once $k$ is chosen.
\end{proof}

The next corollary recovers H{\aa}stad's unconditional hardness theorem
for satisfiable Not-Two and the corresponding three-query PCP consequence
\cite[Theorem~4.3]{Hastad2014NotTwo}, using the earlier reduction of
O'Donnell and Wu.

\begin{corollary}[Satisfiable Not-Two and three-query PCPs]\label{cor:not-two-pcp}
Let $\operatorname{NTW}(a,b,c)$ be the Boolean predicate accepting exactly
zero, one, or three true inputs. For every fixed $0<\eps<3/8$, it is NP-hard
to distinguish a satisfiable instance with constraints
$\operatorname{NTW}(\ell_1,\ell_2,\ell_3)$ from one in which every assignment
satisfies at most a $5/8+\eps$ fraction of the constraints; each literal
$\ell_i$ may be a variable or its negation.
Consequently, every language in NP has a polynomial-time verifier using
$O(\log n)$ random bits and three nonadaptive proof-bit queries, with perfect
completeness and soundness at most $5/8+\eps$, where $n$ is the input length.
\end{corollary}
\begin{proof}
Apply O'Donnell and Wu's reduction
\cite[Theorem~2.1 and Corollary~2.2, author manuscript]{ODonnellWu2008}.
Their premise requires only at-most-$d$-to-1 projections for some fixed $d$.
Take $d=2$, use the uniform distribution on the explicit edge occurrences,
and choose a fixed rational $\delta\in(0,1)$ smaller than their required
soundness threshold. Theorem~\ref{thm:main} supplies this premise with
perfect completeness. Their folded verifier allows all eight literal-sign
variants of Not-Two and gives both conclusions.
\end{proof}

Corollary~\ref{cor:not-two-pcp} also follows by applying the
O'Donnell--Wu reduction to Fei, Minzer and Wang's perfect 4-to-1
theorem~\cite{FeiMinzerWang2026}.

Through Guruswami and Sandeep's reduction at $d=2$
\cite[Theorem~1]{GuruswamiSandeep2020}, Theorem~\ref{thm:main} also
recovers the known NP-hardness of distinguishing three-colorable graphs
from graphs that are not $c$-colorable, for every fixed integer $c\ge3$
\cite[Corollary~1.7]{FeiMinzerWang2026}; the companion
\cite[Theorem~1.1]{OpenAIIndependentSets2026} separately proves stronger
independent-set soundness with three-colorable completeness.
